% !TEX root = ../../../main.tex
\section{Antisymmetric multilinear functions}
\label{sec:linear-algebra-i:antisymmetric-multilinear-functions}

% Sources: October 1 PDF, p. 19 and pp. 25--29; IMG_9210--IMG_9212.
% The PDF uses real vector spaces. The photos use F without a
% characteristic restriction; that broader scope requires author review.
Let \(V\) be an \(n\)-dimensional real vector space.

\subsection{Multilinearity and antisymmetry}

\begin{definition}
  A function \(A:V^n\to\R\) is multilinear if it is linear in each input
  while all other inputs are fixed. Thus, for each position \(i\),
  vectors \(t_1,\ldots,t_{i-1},t_{i+1},\ldots,t_n,v,w\in V\), and
  scalars \(a,b\in\R\),
  \[
    \begin{aligned}
      &A(t_1,\ldots,t_{i-1},av+bw,t_{i+1},\ldots,t_n)\\
      &\quad=aA(t_1,\ldots,t_{i-1},v,t_{i+1},\ldots,t_n)\\
      &\qquad+bA(t_1,\ldots,t_{i-1},w,t_{i+1},\ldots,t_n).
    \end{aligned}
  \]
\end{definition}

\begin{definition}
  A function \(A:V^n\to\R\) is antisymmetric if interchanging any two
  inputs changes its sign. For the transposition \(\tau_{ij}\),
  \[
    A(v_{\tau_{ij}(1)},\ldots,v_{\tau_{ij}(n)})
      =-A(v_1,\ldots,v_n).
  \]
\end{definition}

If two inputs are equal, interchanging them leaves the input unchanged.
Antisymmetry therefore gives
\[
  A(v_1,\ldots,v_n)=-A(v_1,\ldots,v_n),
\]
so \(A(v_1,\ldots,v_n)=0\).

For \(\sigma=\tau_k\circ\cdots\circ\tau_1\), repeated antisymmetry
gives
\[
  \begin{aligned}
    A(v_{\sigma(1)},\ldots,v_{\sigma(n)})
      &=(-1)^kA(v_1,\ldots,v_n)\\
      &=\varepsilon(\sigma)A(v_1,\ldots,v_n).
  \end{aligned}
\]
If \(\phi\in T_n\setminus S_n\), two of the indices
\(\phi(1),\ldots,\phi(n)\) agree, so the corresponding value of \(A\)
is zero. With the extended sign function, both cases are expressed by
\begin{equation}
  A(v_{\phi(1)},\ldots,v_{\phi(n)})
    =\varepsilon(\phi)A(v_1,\ldots,v_n),\qquad \phi\in T_n.
  \label{eq:linear-algebra-i:antisymmetric-reindexing}
\end{equation}

\begin{proposition}
  A multilinear function on a real vector space is antisymmetric if and
  only if it vanishes whenever two inputs are equal. Such a multilinear
  function is also called alternating.
\end{proposition}
% IMG_9211, fact 4, asserts the converse without supplying its proof.
% No polarization argument has been added to complete that proof.

An antisymmetric multilinear function with \(n\) inputs on an
\(n\)-dimensional space is an oriented volume function.

\subsection{Expansion in a basis}

Let \(E=(e_1,\ldots,e_n)\) be a basis of \(V\). Write
\[
  v_j=\sum_{i=1}^n b^i_j e_i,\qquad 1\leq j\leq n.
\]
Expanding the first input by linearity gives
\[
  A(v_1,\ldots,v_n)=\sum_{k_1=1}^n b^{k_1}_1 A(e_{k_1},v_2,\ldots,v_n).
\]
Expanding each remaining input in turn yields
\begin{equation}
  A(v_1,\ldots,v_n)
    =\sum_{k_1=1}^n\cdots\sum_{k_n=1}^n
      b^{k_1}_1\cdots b^{k_n}_n A(e_{k_1},\ldots,e_{k_n}).
  \label{eq:linear-algebra-i:multilinear-basis-expansion}
\end{equation}
Each tuple \((k_1,\ldots,k_n)\) determines a map \(\phi\in T_n\) by
\(\phi(j)=k_j\). The same sum can therefore be written as
\[
  A(v_1,\ldots,v_n)
    =\sum_{\phi\in T_n}b^{\phi(1)}_1\cdots b^{\phi(n)}_n
      A(e_{\phi(1)},\ldots,e_{\phi(n)}).
\]

Now assume \(A\) is also antisymmetric. Terms with repeated indices
vanish, and every remaining index tuple is a permutation of
\((1,\ldots,n)\). Using
\autoref{eq:linear-algebra-i:antisymmetric-reindexing},
\[
  A(e_{\phi(1)},\ldots,e_{\phi(n)})
    =\varepsilon(\phi)A(e_1,\ldots,e_n).
\]
Consequently,
\begin{equation}
  \begin{aligned}
    A(v_1,\ldots,v_n)
      &=A(e_1,\ldots,e_n)
        \sum_{\phi\in T_n}\varepsilon(\phi)
          b^{\phi(1)}_1\cdots b^{\phi(n)}_n\\
      &=A(e_1,\ldots,e_n)
        \sum_{\sigma\in S_n}\varepsilon(\sigma)
          b^{\sigma(1)}_1\cdots b^{\sigma(n)}_n.
  \end{aligned}
  \label{eq:linear-algebra-i:antisymmetric-basis-expansion}
\end{equation}
The second equality uses \(\varepsilon(\phi)=0\) outside \(S_n\).
Thus \(A\) is determined by its value on the ordered basis.

\subsection{Normalization and uniqueness}

\begin{theorem}
  \label{thm:linear-algebra-i:normalized-volume-function}
  For every \(c\in\R\), there is a unique antisymmetric multilinear
  function \(A:V^n\to\R\) satisfying
  \[
    A(e_1,\ldots,e_n)=c.
  \]
\end{theorem}
% PDF p. 29 states existence and uniqueness. The supplied coordinate
% expansion establishes determination by c, but the sources do not verify
% multilinearity and antisymmetry of the formula to prove existence.

In particular, there is a unique such function normalized by
\[
  A(e_1,\ldots,e_n)=1.
\]
If \(B\) is any other antisymmetric multilinear
function with \(n\) inputs, then
\[
  B=cA,\qquad c=B(e_1,\ldots,e_n).
\]
